\section{Testing \& Validation} \label{sec:testingandvalidation}

The telecommunications sector is characterized by stringent \acs{QoS} and performance requirements, inherent to carrier grade networks. With the industry move towards the cloud, it becomes essential to conduct testing procedures to ensure the operation and dimension of cloud-based network elements (i.e. \acs{VNF} and \acs{CNF}) and provide network services that can benefit from cloud technology with slicing and scaling mechanisms. This encompasses laboratory and field testing scenarios, as well as active monitoring and collection of metrics and \acsp{KPI} (e.g. resource utilization) to enable analysis of traffic patterns, relevant for predictive solutions, and closed-loop network automation~\cite{CapacityPerformanceTestMethodsCloudNF}.

\subsection{Data Collection and Analysis via NWDAF}

An important area of research corresponds to the development of \acs{ML} \acs{NWDAF} models for predicting \acs{DDoS} attacks or implementing autoscaling tecniques to dynamically adjust network resources in response to high traffic conditions. \citep{UserTerminalsAsAttackers, AdvancingPredictiveSecurity} utilize an annotated dataset to detect \acs{DDoS} attacks collected from a \acs{5G} testbed. The studies focuses mainly on the validation of the \acs{ML} algorithms and analysis of the traffic profile of the benign (using Skype and Youtube) and malicious \acsp{UE} (performing the various \acs{DDoS} attack types). Despite specifying that there are a total of nine \acsp{UE} (five of them malicious), each attributed to a specific cell, the underlying functional implementation of traffic generation and data collection are not specified. However, the set of features collected, from the various \acsp{NF}, and utilized by the algorithm constitute key insights for such prediction endeavors. \citet{PerformanceComparisonNWDAFSecurity} address the same use case, distinguishing between Control Plane and Data Plane \acs{DoS} scenarios. In the first scenario, anomalous behavior was characterized by very quick and frequent association-deassociation of \acsp{UE}, to flood the network control plane, whereas the second scenario performed the \acs{DoS} attack utilizing the hping tool deployed in the \acs{UE}. The data collection process is implemented through a pipeline that periodically issues asynchronous, non-blocking requests to WebSocket \acsp{API}, provided by the \acs{5G} platform powered by Amarisoft Callbox Classic. The collected metrics include downlink and uplink bitrates, transmission and error counts, session management statistics and \acs{UE} registration records. Nonetheless, the study lacks validation procedures to evaluate the correctness of the data collection process under real-world conditions.

In contrast, \cite{DataCollectionUsingNWDAF5GCN} validates a custom \acs{NWDAF} implementation and the coherence of the metrics collected (i.e. comparison of application-level metrics versus metrics collected by the \acs{NWDAF}), according to generated traffic, within a \acs{5G} system in a Linux-based \acs{VM} testbed. The testbed is constituted of a \acs{RAN} \acs{VM} connecting a \acs{UE} \acs{VM}, containing all the \acsp{UE}, to a intermediary \acs{UPF} (I-\acs{UPF}), which in turn connects to an anchor \acs{UPF} (A-\acs{UPF}) that hosts the server-daemon pairs for the \acsp{UE}. The \acs{NWDAF} metrics are obtained at the I-\acs{UPF} for uplink data volume and \acs{RTT}, while the A-\acs{UPF} measures the downlink data volume. In order to generate traffic, a custom traffic loader was developed in Go and Bash, using iPerf3 v10 to produce synthetic traffic based on predefined \acs{YAML} configuration files with configurable uplink and downlink traffic models for \acs{TCP} or \acs{UDP}. This tool is composed by two components: a simulated \acs{UE} and a corresponding server-side daemon that receives (uplink) or sends (downlink) the requested byte stream. By combining it with isolated network namespaces with different \acsp{IP}, it's possible to synthesize traffic with \acs{PDU} session granularity, between different source \acsp{IP} and the network, representing the \acsp{UE} served by a network slice. With this setup, four different traffic profiles are configured to test the \acs{NWDAF}: \acs{TCP}-Long, \acs{TCP}-Short, \acs{UDP}-\acs{CBR}, \acs{UDP}-\acs{VBR} and a combined profile incorporating them all. The average bit rate (throughput) of these traffic profiles is then compared to the \acs{NWDAF}'s, corroborating the collected metrics. To validate latency measurements, the \acs{NWDAF} statistics were compared against the ping result among the \acsp{UPF} serving a slice with variable background conditions. In this case, the traffic generation process was achieved manually, for the sake of validating the collection of data by the \acs{NWDAF}. However, ideally this would be achieved automatically without human intervention, through continous validation of the \acs{NWDAF} metrics in real time or, at least, through the use of dedicated validation platforms that can automate traffic generation and other related procedures.

Despite detailing the collection of metrics, the previous \acs{NWDAF} implementation is limited to the \acs{UPF} and operates at a low level within a Linux-based \acs{VM} testbed, not realistic for real world environments. By comparison, \cite{ExplorationNWDAFDevelopmentArchitecture6G} presents a full architecture for the \acs{NWDAF} with data collection and management, and model training and management components, for a vertical scaling use case within a testbed composed by a Open5GS Kubernetes \acs{CN} (with \acs{AMF}, \acs{SMF} and \acs{UPF}) and \acs{OAI} integrated with \acs{USRP} hardware. Two \acsp{UE} are connected to the \acs{CN}, with traffic load being generated by transmitting data packets between them using the iPerf tool. During the experiment, real-time data is collected via Prometheus from the \acs{AMF}, \acs{SMF} and \acs{UPF} to monitor the experiment and trigger \acs{UPF} vertical scaling actions based on \acs{UPF} \acs{CPU} usage. The analysis of this data corroborated the increase of \acs{AMF}/\acs{SMF} traffic upon \acs{UE} connection and of the number of N4 sessions due to data transmission. Lastly, the \acs{UPF} traffic values remained at a constant level showcasing  successful communication between devices, with intensive data transmission triggering \acs{UPF} scaling across three different \acs{ML} models with a lower response time than a 70\% \acs{CPU} threshold based strategy.


\subsection{Network Applications} \label{sec:testingandvalidation:netapps}

The development of Network Applications needs to be accompanied by testing procedures, leveraging monitoring stacks and collected metrics to validate usage scenarios in real testbeds. This section reviews the testing methodologies and validation mechanisms reported in the literature, focusing on how they're used to assess and demonstrate the feasibility of some of the use cases discussed throughout this document.

In the Section \ref{sec:testbeds:5GC:5GC}, it was presented a dynamic \acs{QoS} adaptation scenario based on the CAMARA Device Location and \acs{QoD} \acsp{API}, in the context of a vehicular remote monitoring digital twin use case~\cite{DynamicQoSAdaptationviaExposedServiceAPIs}. A testing experiment is performed to showcase the dynamic \acs{QoS} modifications at a designated geographical area, represented by a cell ID. As only a single \acs{gNB} is available, radio cell changes are emulated by periodically modifying the cell ID of connected \acsp{UE} at the \acs{AMF}, simulating movement between different geographical areas. Two \acsp{UE} continously generate iperf \acs{TCP} uplink traffic, at a combined data traffic of 10 Mbit/s. The \acs{V2X} service subscribes to location updates for one of the \acsp{UE} and upon entering a specific area (the north cell) it enforces a \acs{QoS} profile (via the CAMARA \acs{QoD} \acs{API}) with absolute priority of the selected \acs{UE}'s data flow over other traffic, automatically removing it when the \acs{UE} leaves. This process is performed 200 times to collect statistically relevant measurements. The data collection and analysis is achieved through the use of Prometheus and Grafana, at the level of the server infrastructure (OpenStack). In result, the \acsp{UE} compete equally in the absence of any high-priority \acs{QoS} profile, with each \acs{UE} generating roughly 5 Mbit/s of throughput, as opposed to the selected \acs{UE} displaying 9 Mbit/s of throughput when the \acs{QoS} profile is enforced on the north cell. The latency derived of the combined usage of both \acsp{API} was also measured between the network detecting a location change (with a \acs{NGAP} location report at the \acs{AMF}) and the \acs{QoS} profile being confimed at the \acs{gNB}, resulting in 50th percentile latency of 255.0 ms, 90th percentile of 282.8 ms, 95th percentile of 290.3 ms, and 99th percentile of 337.0 ms. These latency values essentially represent an order of magnitude of the overall latency, being dependant on the specific implementation and geographical area, and could be further reduced by employing predictive models that antecipate when the \acs{QoS} profile shoud be applied.

A similar approach is followed in a tele-operated driving context, leveraging the same \acs{QoD} \acs{API} and a performance metrics \acs{API} to provide subscribers with periodic performance data for a device. The tele-operation service is assigned the high priority \acs{QoS} profile, depending on the test case, whereas the general-purpose default \acs{QoS} profile is assigned to a \ac{UE} that continuously sends background iperf TCP traffic up to the maximum achievable uplink throughput of roughly 30 Mbit/s. Similarly to the previous study, when the \acs{QoD} was disabled the \acs{UE} competed for the bandwidth reaching an higher throughput (15 Mbit/s compared to 10 Mbit/s achieved by the vehicle) and resulting in a lower average frame rate of 15 \acs{FPS} and an \acs{E2E} network latency of up to 1200 ms. Setting the high priority \acs{QoS} profile (manually through a GUI), stabilized the \acs{FPS} and \acs{E2E} latency at substantially lower values, as the data traffic coming from the vehicle has priority over the background traffic. The metrics collection process relies on an \acs{NWDAF} implementation leveraging the testbed Prometheus exporter for retrieving \acs{UPF} data, with the \acs{E2E} latency corresponding to the total delay between a network packet being sent by the vehicle and received at the remote station.

Both previous approaches validate their use cases demonstrating how throughput varies under the applied \acs{QoS}. However, this validation approach is inherently limited to use cases focused on \acs{QoS} management. \cite{VirtualOnBoardUnitMigration} deviates from this use case and validates the automatic migration of a v\acs{OBU} between \acs{MEC} domains as a vehicle moves between \acs{5G} cells. The implementation was tested in two different parts: first in a controlled indoor lab the \acs{OBU} interacts with a \acs{NEF} Emulator to subscribe to cell change events triggered by a simulated device on a predefined path, followed by an outdoor experiment where an actual vehicle (with an \acs{OBU}) travels between coverage areas to trigger a physical handover in an outdoor \acs{5G} network. Openstack was the virtual infrastructure manager used, with \acs{OSM} being chosen as the orchestrator to manage the dynamic deployment of the v\acs{OBU} \acs{5G} Network Application. In both test cases, the v\acs{OBU} migration to a new \acs{MEC} domain was successful, with the latter being corroborated by a spike on \acs{CPU} usage of the Openstack compute node located in the given \acs{MEC} domain the migration occured to.


\subsubsection{Testing Pipelines} \label{sec:pipelines}

Despite validating the basic functionality of the implementations for each case, there is still insufficient adherence to testing pipelines with automated processes for provisioning each of the components present in the respective validation scenario. Instead most testing procedures still rely on low-level scripts or manual workflows to evaluate their applications. Therefore, in this section it's performed an overview of the various validation procedures, with different levels of automation, present in the literature.

\citet{ExperiencesChallengesTestingValidatingNetworkAware} explore the validation of network-aware \acs{PPDR} \acs{XR} applications in the Patras\acs{5G} testbed. To this end, \acs{OSL} was used as the \acs{OSS}, being responsible for provisioning and orchestrating the depoyment of all components to provide a private \acs{5G} network for testing. This includes any appropriate service that's available through the \acs{OSL} service catalog, including \acs{5G} network capabilities and \acsp{gNB}. Telemetry and monitoring tools are integrated in the Patras\acs{5G} testbed by using Prometheus, providing metrics about the \acs{RAN}, along with additional metrics obtained via custom data collectors. Then, for each of the three tested \acfp{UC}, \acsp{VM} were provisioned according to their specific requirements (specs, network access, traffic requirements, \acs{GPU} access). Such \acsp{VM} also include simulated \acsp{UE}. During the lab trials of each \acs{UC} metrics for bandwidth consumption and application latency were measured. The testing methodology combined remote testing for \acs{UC} owners to install and configure their applications on the \acsp{VM}, while testbed personnel managed the network infrastructure and any physical devices, followed by physical testing mostly related to including physical \acs{AR}/\acs{VR} equipment in the testing process. Consequently, despite usign \acs{OSL} for orchestration and provisioning, the process remains partially automated, requiring a certain level of manual configuration for tasks such as the installation and setup of applications in the provisioned \acsp{VM}.

There are more complete solutions within the literature, including a full framework to orchestrate, test, and validate network applications based on the 5GASP project~\cite{AutomatedTestingOrchestrationFramework}. This framework addresses the full workflow of (i) onboarding the network application, (ii) their orchestration, and (iii) the orchestration of the tests to be performed to the network application. The network application and its testing processes are packaged as a bundled service within \acs{OSL}, which also includes a Network Slice Template (based on \acs{GSMA}'s \acs{GST}) to allow testing the application under specific network conditions. For orchestrating the network application, \acs{OSL} interacts with \acs{OSM} via the interface SOL005, supporting the deployment of network applications in both Openstack and Kubernetes. The test cases are defined by a 5GASP Testing Descriptor, in \acs{YAML}, and carried out via the Robot Framework, which automatically produces report and log files in \acs{HTML}. Based on this blueprint, a \acs{CI} Service Testing Orchestrator orchestrates the testing processes, dynamically generating pipeline configuration files and submitting them to Jenkins \acfp{TEA} for execution. After collecting and running all test cases, the \acsp{TEA} aggregate the testing outputs and forward them to the \acs{CI} Service, which creates a visual testing reporting on \acs{OSL} for testers. Regarding monitoring, the computing and resource usages of network app components are scraped from Openstack and Kubernetes using Prometheus, and visualized via Grafana, with measurements on the \acs{5G} resources consumed by each network app being obtained through a man-in-the-middle monitoring switch between the \acs{NFVI}/\acs{CNI} and the \acs{5G} system. This switch handles all interaction between the network applications and the \acs{5G} \acsp{UE}, allowing extensive \acs{KPI} validation and monitoring of throughput, \acs{RTT}, jitter and packet loss percentage. The testing framework was applied in the validation of 11 different network applications across 4 different agnostic testing scopes: \acs{5G} readiness, Security \& Privacy, Performance \& Scalability and Availability \& Continuity.

Although offering a comprehensive DevSecOps testing solution, the framework does not follow a phased methodology that ensures each deployed artifact operates as expected before progressing to deploy additional artifacts. In this context, the following approach extends on the previous work by designing a multistage testing pipeline to support the design, orchestration, and lifecycle management of \acs{5G} resources, such as network slices, and applications~\cite{TMFOpenDigitalArchitecture5GServiceOrchestration}. Similarly to the previous approach, the \acs{OSL} corresponds to the entrypoint of the framework, providing \acs{TMF}-aligned data models and \acsp{API}. \acs{GSMA}'s \acsp{NEST} are mapped into \acs{OSL} services to offer configurable network slices tailored to specific use cases. Once ordered from the \acs{OSL} Service inventory, it is converted  into actionable data models to be consumed by an orchestrator capable of provisioning the network slice, which then becomes accessible in \acs{OSL} for full lifecycle management. This approach is also adopted for the provisioning of other \acs{5G} resources, such as an \acs{UE}. Regarding orchestration, the \acs{OSL} provides interfaces for orchestrating \acs{NFV} (using \acs{OSM} through \acs{ETSI} SOL005) and containerized (using kubernetes through CRIDGE component) applications. To offer monitoring components and simplified CAMARA \acsp{API} for the management of 5G resources and applications, \acs{TMF} services are again utilized, deploying Prometheus/Grafana monitoring bundles via \acs{OSM} and CAMARA \acsp{API} through Kubernetes.

For performing the testing procedure, the previous DevOps testing solution~\cite{AutomatedTestingOrchestrationFramework} is integrated via the \acs{TMF} 653 \acs{API} (Service Test Management) to support a multi-step approach with the following steps: (i) onboard/design/configuration of \acs{5G} resources and application, (ii) order network slice provisioning, (iii) order \acs{5G} \acs{UE} provisioning, (iv) validate network slice (as per the \acs{SLA}), (v) order application, (vi) validate application, (vii) order monitoring stack, (viii) order CAMARA \acsp{API}, with the service becoming ready for use. The provisioning of network slices and \acsp{UE} was achieved by using a proprietary network slicing \acs{API} from 5GAIner. A computer vision application for monitoring an assembly line, which leverages a CAMARA \acs{QoD} Provisioning \acs{API} to support diverse video streaming requirements, was tested under this pipeline. The first validation step evaluates the network slice capabilities using the provisioned \acs{UE} in regards to throughput, latency, packet loss, jitter, and power consumption, with iPerf3, ping, and a custom Python script being used to measure the latency between the \acs{5G} \acs{UE} and the application. The second validation stage executed 2 custom test cases using a healthcheck to verify that the application's \acs{AI} model is running and it can consume video streams from the \acs{5G} \acs{UE}. 

\hide{
\later{se tiver tempo (não vou ter xD): NSSF Function in 6G Networks Based on MLOps Deployment Model}
}

\subsection{Metrics \& KPIs}
\label{sec:testingandvalidation:metrics}

In the preceding sections, the discussion was focused on the various approaches for data collection and analysis via \acs{NWDAF}, as well as the validation of network applications, with special emphasis on the automation and orchestration of testing pipelines. Across all these contexts, the systematic collection of metrics plays a central role in ensuring the correctness, efficiency and reliability of the deployed components. Such metrics, even those from \acs{NWDAF} scenarios, can provide key insights into the operational behavior of network elements and applications.

Within the context of this dissertation, these metrics are essential for the design and validation of \acs{5G} network applications, including the required monitoring solutions. To provide a consolidated view of the most relevant metrics discussed throughout the literature, Table \ref{tab:metrics:all} summarizes the metrics, attributes \& \acsp{KPI} collected across all the aforementioned scenarios.

\begin{table}[ht]
\centering
\tiny
\renewcommand{\arraystretch}{1.3} % Increase row height
\rowcolors{2}{gray!15}{white}     % Alternating row colors
\begin{tabular}{p{4cm} p{11cm}}
\toprule
\textbf{Source} & \textbf{Metrics, attributes \& \acsp{KPI}} \\
\midrule
Advancing Predictive Security~\cite{AdvancingPredictiveSecurity} & \acf{UL} and \acf{DL} (throughput, \acs{CRC}, \acs{MCS}, \acs{PDCP} \acs{SDU}), \acs{QoS} flow identifier, Slice service type, \acs{UE} aggregate maximum bitrate (\acs{UL} and \acs{DL}), \acs{UE} transmission power, \acs{CQI} \\
Exploration of \acs{NWDAF} Development Architecture~\cite{ExplorationNWDAFDevelopmentArchitecture6G} & \acs{UPF} \acs{CPU} usage, data transfer rate, \acs{UE} \acs{UL}/\acs{DL} troughput, N4 session counts \\
Network Slice Selector in \acs{IoT} Services~\cite{NetworkSliceSelectorIoT} & Latency, jitter, packet loss, download/upload rates, distance, reliability, user density \\
Cloud-Enabled Deployment of \acs{5G} Core Network~\cite{CloudEnabledDeployment5GN} & \acs{E2E} latency, throughput, total aggregated maximum bitrate (\acs{UL}/\acs{DL}) \\
Dynamic QoS Adaptation Via Exposed Service \acsp{API}~\cite{DynamicQoSAdaptationviaExposedServiceAPIs} & Throughput (uplink/downlink), cell id, latency \\
Load Balancing in Connected Vehicle Services~\cite{LoadBalancingConnectedVehicleServices} & average response times, \acs{CPU} utilization, \acs{RTT} \\
An Automated Testing and Orchestration Framework~\cite{AutomatedTestingOrchestrationFramework} & Response time, throughput, \acs{RTT}, jitter, packet loss \\
Virtual On-Board Unit Migration~\cite{VirtualOnBoardUnitMigration} & latency, \acs{CPU} usage \\
Next-Gen Connectivity: Dynamic QoS Optimization~\cite{NextGenConnectivityDynamicQoS} & Round-trip latency, upload/download speeds, \acs{QoS} profile adaptation time \\
On the benefits and caveats of exploiting Quality on Demand~\cite{BenefitsCaveatsQoDNetworkAPIsVideoStreaming} & Stall duration, average video bitrate, frequency/duration of bandwidth boosts \\
On-Demand Network QoS Adaptation~\cite{OnDemandNetworkQoSAdaptation} & Uplink/downlink rates (throughput), \acs{E2E} network latency, video quality (\acs{FPS}) \\
Design and Evaluation of a \acs{NDT}~\cite{DesignAndEvaluationNetworkDigitalTwin} & \acs{UL}/\acs{DL} throughput \\
Performance Comparison of NWDAF-Based Security~\cite{PerformanceComparisonNWDAFSecurity} & \acs{UL}/\acs{DL} bitrates, transmission/error counts \\
DEMO: On-demand 5G/6G edge verticals~\cite{DemoOnDemand5GEdgeVerticalsViaThirdPartyUPF} & \acs{E2E} latency, throughput, \acs{RTT} \\
Capacity and Performance Test Methods~\cite{CapacityPerformanceTestMethodsCloudNF} & \acs{CPU}, memory, storage, network utilization; call success rate; data throughput \\
Data Collection Using \acs{NWDAF}~\cite{DataCollectionUsingNWDAF5GCN} & Uplink/downlink traffic volume on \acs{TCP} and \acs{UDP}, \acs{RTT}, average bandwidth \\
Experiences and challenges on testing and validating Network Aware \acs{PPDR} \acs{XR}~\cite{ExperiencesChallengesTestingValidatingNetworkAware} & \acs{UL}, application latency, video image quality (\acs{FPS}) \\
Pushing the Boundaries of Scalable 5G Core~\cite{pushingBoundariesScalable5GCN} & Latency, throughput, response time, resource usage (\acs{CPU}\% and \acs{RAM}) \\
Optimizing \acs{3GPP} \acs{5G} \acs{NEF} Integration~\cite{optimizing5GNEFIntegrationThroughNEFandUPF} & latency, \acs{QoS} adjustment speed, resource allocation efficiency \\
Capability Exposure Vitalizes \acs{5G}~\cite{CapabilityExposureVitalizes5GNetwork} & \acs{RTT}, effective time of routing policy (uplink classifier insertion completion), downlink throughput \\
\acs{TMF} Open Digital Architecture for \acs{5G} Service Orchest.~\cite{TMFOpenDigitalArchitecture5GServiceOrchestration} & throughput (\acs{UL}/\acs{DL}), latency, packet loss, jitter, \acs{UE} energy consumption, network slice, \acs{UE} and monitoring provisioning times \\
\bottomrule
\end{tabular}
\caption{Collected Metrics, attributes and \acsp{KPI} in \acs{5G} networks}
\label{tab:metrics:all}
\end{table}